\documentclass[12pt]{article}
\usepackage[T1]{fontenc}
\usepackage[utf8]{inputenc}
\usepackage{lmodern}
\usepackage{textcomp}
\usepackage{enumitem}
\usepackage{geometry}
\usepackage{graphicx}
\usepackage{amsmath, amssymb}
\geometry{margin=1in}
\begin{document}
\begin{center}
Political Science - Graduate Level Assessment 2 \
Total Marks: 40 \
Answer all questions.
\end{center}

\section*{Multiple Choice (10 marks)}
\begin{enumerate}[label=\arabic*.]
\item What fact impinged upon America's 'neutral' stance towards the belligerent in World War I between 1915 and 1917?
\begin{enumerate}[label=(\alph*)]
    \item US banks funded the allies much more than Germany and Austria-Hungary
    \item US banks funded Germany and Austria-Hungary much more than the allies
    \item Wilson secretly favoured the German position
    \item None of the above
\end{enumerate}
\item Which of the following is the odd one out?
\begin{enumerate}[label=(\alph*)]
    \item Institutional constraints and resource constraints.
    \item Information and time horizon.
    \item Goals.
    \item Vectors of economic incentives.
\end{enumerate}
\item Why did George H.W. Bush allow Saddam Hussein to remain in power after the Gulf War of 1991?
\begin{enumerate}[label=(\alph*)]
    \item Lack of US firepower
    \item Concern over oil supplies
    \item Limited UN mandate and fear of a protracted conflict
    \item Difficult terrain and fear of civilian casualties
\end{enumerate}
\item What was meant by the term 'American multiplication table'?
\begin{enumerate}[label=(\alph*)]
    \item Increase in the US population
    \item Increase in US finances
    \item Increase in US military capability
    \item Increase in US international influence
\end{enumerate}
\item What features distinguish Socio-Economic accounts of US Cold War foreign policy?
\begin{enumerate}[label=(\alph*)]
    \item A focus on class and economic interests
    \item A focus on liberal ideology
    \item A focus on language and culture
    \item All of the above
\end{enumerate}
\end{enumerate}

\section*{True / False (10 marks)}
\textit{Answer True or False}
\begin{enumerate}[label=\arabic*.]
\setcounter{enumi}{5}
\item True or False: Societal groups are formed through complex social interactions among individuals within a community.
\item True or False: Rogue States are defined as nations that operate outside the recognized norms of international relations and diplomacy.
\item True or False: NSC-68 expanded U.S. strategy by globalizing and militarizing containment while advocating for the development of the hydrogen bomb.
\item True or False: George H.W. Bush maintained a cautious support for both Mikhail Gorbachev and Boris Yeltsin during their leadership in Russia.
\item True or False: Constraints on US foreign policy decision making include the foreign policies of other states, international law, and intergovernmental organizations.
\end{enumerate}

\section*{Long Form Response (20 marks)}
\textit{Answer each question with a clear and complete explanation.}
\begin{enumerate}[label=\arabic*.]
\setcounter{enumi}{10}
\item Discuss the role of various economic tools in security policy and critically analyze why diplomacy is considered distinct from these tools in shaping international relations.
\item Discuss the key principles of the Responsibility to Protect (R$_{2}$P) doctrine and critically analyze why the principle of sovereignty is often seen as a limitation rather than an element of R$_{2}$P.
\end{enumerate}

\vfill
\noindent\textit{End of Paper}
\end{document}